\subsection{The prescribed spectral curve and BA reconstruction}
\label{sec:BA-spectral-converse}

The interpolation functions $h_L,h_M$ give spectral coordinates on
the prescribed curve. With several punctures, identifying that curve
with the normalized characteristic curve requires the birationality
condition stated below.

\begin{proposition}[The spectral curve of a prescribed realization]
\label{prop:BA-prescribed-spectral-curve}
For the phase coefficients \cref{eq:general-phase-values}, put
\begin{equation}\label{eq:BA-prescribed-spectral-map}
 \nu(P)=-\tfrac12(h_L(P)+h_M(P)),\qquad
 \lambda(P)=\tfrac12(h_M(P)-h_L(P)).
\end{equation}
The characteristic polynomial of the canonical pencil vanishes on
this map. If the map is birational onto its image, $\mathcal C$
normalizes the full characteristic curve. A sufficient condition is
that one puncture has a finite limiting $\lambda$ distinct from the
other punctures, in particular pairwise distinct $b_k^N$.
For $N=1$ this condition is automatic.

Let $\Gamma$ be the projective plane characteristic curve. Write
$d=L+M+N$ and $G=LM+MN+NL-d+1$. On this birational locus
the total singularity defect beyond the ordinary three infinity
multiple points is
\begin{equation}\label{eq:BA-prescribed-genus-defect}
 \sum_{q\in\operatorname{Sing}(\Gamma)}\delta_q
 -\binom L2-\binom M2-\binom N2=G-g.
\end{equation}
For distinct phases within groups the additional defect is affine.
\end{proposition}
\begin{proof}
Let $\mathsf X_{\rm BA}$ be the block matrix in
\cref{eq:canonical-pencil} with the raw arrays
\cref{eq:multipoint-mode-arrays} and the phase diagonals
\cref{eq:general-phase-values}. The explicit vector is
\begin{equation}\label{eq:BA-prescribed-eigenvector}
 r(P)=\begin{pmatrix}
 (-\kappa_i\Omega(Q_i,P)/c_i)_{i=1}^{L}\\
 (c_{L+j}\Omega(Q_{L+j},P))_{j=1}^{M}\\
 (e_k\psi_k^\vee(P))_{k=1}^{N}
 \end{pmatrix}.
\end{equation}
The first two block rows of
$(\mathsf X_{\rm BA}+\nu(P)\mathsf D)r=\lambda(P)r$
are the convolution identities for $h_M,h_L$. The third follows by
subtracting their dual weighted equations:
\[
 \sum_{i=1}^{L}\kappa_i\psi_k^\vee(Q_i)\Omega(Q_i,P)
 -\sum_{j=L+1}^{K}\kappa_j\psi_k^\vee(Q_j)\Omega(Q_j,P)
 =(\lambda(P)-b_k^N)\psi_k^\vee(P).
\]
The phase transformation \cref{eq:multipoint-phase-gauge} conjugates
the pencil and vector by the same diagonal matrix, preserving this
spectral equation. Adding a scalar origin $\rho\Id_d$ to the
pencil shifts $\lambda$ by $\rho$ and preserves the eigenvectors.

The poles of $\nu$ are simple, at all $Q_a,R_k$. Their branches obey
\[
 \lambda-\nu\to-2b_i^L\ (P\to Q_i),\qquad
 \lambda+\nu\to-2b_j^M\ (P\to Q_{L+j}),\qquad
 \lambda\to b_k^N\ (P\to R_k).
\]
If one $R_k$ has a unique finite limiting ordinate, it is the sole
point over its image in the normalization. Its simple pole of $\nu$
makes this point unramified, so the covering degree is one.
The same argument applies to a uniquely separated outer branch.
For a birational map, $\deg\nu=d$ implies that the minimal polynomial
of $\lambda$ over $\C(\nu)$ has degree $d$. The monic characteristic
polynomial has that degree and vanishes, so it is this minimal
polynomial. The plane degree is $d$; subtracting the ordinary
infinity defects gives \cref{eq:BA-prescribed-genus-defect}.
\end{proof}

\paragraph{Prescribed curves and spectral fibres.}
Fix the marked spectral curve and its functions $\lambda,\nu$,
with distinct phase intercepts. On the regular framed spectral
correspondence, a smooth genus-$G$ level has $G$ divisor-class
parameters and $d-1$ independent mode gauges, hence dimension
$G+d-1=p$. A prescribed genus-$g$ normalization instead supplies
$g+d-1$ parameters after identifying equivalent meromorphic frames.
Relative to an expected-dimensional spectral fibre, this realized
family has codimension $G-g$. It is an invariant normalization-based
family on a special level when $g<G$, and is open in the regular
fibre when $g=G$.

This family has a concrete spectral-sheaf interpretation. Separate
the forced ordinary infinity points, and let
$n:\mathcal C\to\Gamma_{\mathrm{sp}}$ normalize the remaining
integral spectral model. At the reference time the components of
\cref{eq:BA-prescribed-eigenvector} are sections of
\[
 \mathscr L=\mathcal O\left(\mathscr D^\vee+
                         \sum_{a=1}^{K}Q_a\right).
\]
Subtracting the full pole fibre of $\nu$ gives
$\mathscr L(-\nu^{-1}(\infty))=\mathcal O(\mathscr D^\vee-
\sum_{k=1}^{N}R_k)$, with both cohomology groups zero.
Its leading-value sequence shows that the $d$ components form a
basis and that $\nu_*\mathscr L\simeq\mathcal O^d$.
Multiplication by $\lambda$ gives the pencil, so the component
spectral sheaf on $\Gamma_{\mathrm{sp}}$ is $n_*\mathscr L$.
Thus the realized family is the stratum obtained by pushing line
bundles forward from the normalization, with the mode framings
retained. Its divisor-class dimension is $g$.
The full compactified spectral fibre can also contain sheaves with
singularity gluing data \cite{AdamsHarnadHurtubise1990II,Beauville1990}.
The general solution of
\cref{thm:BA-spectral-converse} reconstructs the curve belonging to
generic initial modes. For repeated phases its active dimensions,
moment data and restored block bases are accounted for separately
in Appendix~\ref{app:BA-active-converse}.

With several punctures, the interpolation functions can instead factor
through a covering; curve identification then requires birationality.
The stronger completion in \cref{cor:BA-finite-list} permits distinct
differences $a_{M,k}-a_{L,k}$ to be prescribed, giving the sufficient
branch separation used in the proposition.
For $N=1$, the product coordinates $\lambda-\nu,\lambda+\nu$ have
degrees $M+1,L+1$; their image has bidegree $(L+1,M+1)$ and arithmetic
genus $LM$, recovering the one-point defect $LM-g$.

\begin{theorem}[Generic BA reconstruction]\label{thm:BA-spectral-converse}
Fix positive $L,M,N$ and arbitrary phase coefficients. Generic
canonical initial modes reconstruct scalar multipoint BA bases and
their kernel on a smooth normalized active spectral curve.
For distinct phases within groups all directions are active and the
genus is $G=LM+MN+NL-L-M-N+1$; the last group supplies $N$
punctures, and the ordinary and dual divisors have degree $G+N-1$.
For repeated phases, use constant canonical bases within equal-phase
blocks and remove inactive directions as in
Appendix~\ref{app:BA-active-converse}. Apply the construction with
active dimensions $L_a,M_a,N_a$ and genus
\cref{eq:BA-active-genus}, then restore the zero directions and bases.
The original coordinates can then be constant linear combinations
of the active scalar BA evaluations.
\end{theorem}
\begin{proof}
Work in the $H_0$ phase gauge of \cref{eq:canonical-phase-gauge},
omitting its tildes during this proof. The pencil retains the diagonal
blocks $-2\mathsf B_L,-2\mathsf B_M,\mathsf B_N$, and its time matrix is
\begin{equation}\label{eq:multipoint-H0-time}
 \mathsf T_0(\nu)=
 \begin{pmatrix}0&0&\bar{\boldsymbol w}^{\trans}\\
 0&0&-\boldsymbol v\\
 -\boldsymbol w&\bar{\boldsymbol v}^{\trans}&-\nu\Id_N
 \end{pmatrix}.
\end{equation}
For distinct phases the regular normalized curve has branches
\begin{equation}\label{eq:BA-converse-infinity}
\begin{aligned}
 \lambda&=\nu-2b_i^L+O(\nu^{-1})&&\text{at }Q_i,\quad1\leq i\leq L,\\
 \lambda&=-\nu-2b_j^M+O(\nu^{-1})&&\text{at }Q_{L+j},\quad1\leq j\leq M,\\
 \lambda&=b_k^N+O(\nu^{-1})&&\text{at }R_k,\quad1\leq k\leq N.
\end{aligned}
\end{equation}
The nonempty regular locus, including repeated-phase active
normalizations, is proved in \cref{prop:BA-active-regularity}.
Put $s_k^{-1}=-\nu$ at $R_k$. The functions
\begin{equation}\label{eq:BA-converse-h}
 h_L=-\lambda-\nu,\qquad h_M=\lambda-\nu
\end{equation}
have the required poles and satisfy
$h_M(Q_i)=-2b_i^L$, $h_L(Q_{L+j})=2b_j^M$,
$a_{L,k}=-b_k^N$, $a_{M,k}=b_k^N$.
The interpolation and phase conditions thus come from the initial
spectral curve, rather than restricting initial modes further.

Let $f_k$ be the last $N$ coordinate vectors, $\mathsf J$ their
coordinate projection, and $a=\sum_{k=1}^{N}f_k$.
Choose right and left spectral eigenvectors $r,\ell$ with
$i\dot r=-\mathsf T_0r$, $i\dot\ell=\ell\mathsf T_0$ and reference
normalizations $a^{\trans}r(P,0)=\ell(P,0)a=1$. Define
\begin{equation}\label{eq:BA-converse-functions}
 \psi_k(P)=\ell(P)f_k,\qquad
 \psi_k^\vee(P)=f_k^{\trans}r(P),\qquad
 \Omega(P,Q)=-\frac{\ell(P)r(Q)}{\nu(P)-\nu(Q)}.
\end{equation}
Since $\mathsf T_0(\nu(P))-\mathsf T_0(\nu(Q))
=-(\nu(P)-\nu(Q))\mathsf J$, this kernel satisfies
\cref{eq:general-kernel-rate}. The eigenvalue equations remove apparent
poles at distinct points of an equal-$\nu$ fiber. At $R_k$, the leading
eigenvectors are its coordinate direction and the temporal equation
gives exactly \cref{eq:general-BA-asymptotics}.

For the divisor count use the positive-degree spectral component
line bundle $\mathscr L$, whose pushforward under $\nu$ is
$\mathcal O^d$ \cite{Beauville1990}. Its degree is $G+d-1$.
The last $N$ components vanish simply at the $L+M$ outer infinity
points. Their remaining common BA divisor has degree $G+N-1$.
More explicitly, put $\mathscr E=\mathscr L(-\nu^{-1}(\infty))$.
Then $H^0(\mathscr E)=H^1(\mathscr E)=0$, since its pushforward is
$\mathcal O(-1)^d$. Adding the $N$ punctures gives an exact
leading-value isomorphism
$H^0(\mathscr E(\sum_{k=1}^{N}R_k))\simeq\C^N$.
This proves the required BA uniqueness. The transpose module gives
the dual statement; the analytic infinity justification for repeated
phases is given in Appendix~\ref{app:BA-active-converse}.
The duality differential is
\begin{equation}\label{eq:BA-converse-duality}
 d\Omega_0=-\frac{d\nu}{\ell(P)r(P)}.
\end{equation}
It is regular at the outer branches, has the required double poles
at $R_k$, and its apparent projection-branch poles cancel with $d\nu$.

The reference normalization gives a concrete recovery chart. Set
\begin{equation}\label{eq:BA-converse-constants}
\begin{aligned}
 c_i&=\sum_{k=1}^{N}w_{ki}(0),&
 \bar c_i&=\sum_{k=1}^{N}\bar w_{ki}(0),&
 \kappa_i&=c_i\bar c_i,\quad1\leq i\leq L,\\
 c_{L+j}&=\sum_{k=1}^{N}v_{jk}(0),&
 \bar c_{L+j}&=\sum_{k=1}^{N}\bar v_{jk}(0),&
 \kappa_{L+j}&=c_{L+j}\bar c_{L+j},\quad1\leq j\leq M.
\end{aligned}
\end{equation}
The barred constants are independent partner normalizations, not
complex conjugates; the weights are products of initial sums.
On the open chart where these sums are nonzero, outer eigenvector
expansions give, for $1\leq i\leq L$, $1\leq j\leq M$ and $1\leq k\leq N$,
\begin{equation}\label{eq:BA-converse-recovery}
\begin{aligned}
 \psi_k^\vee(Q_i)&=w_{ki}/c_i,&
 \psi_k(Q_i)&=\bar w_{ki}/\bar c_i,\\
 \psi_k(Q_{L+j})&=v_{jk}/c_{L+j},&
 \psi_k^\vee(Q_{L+j})&=\bar v_{jk}/\bar c_{L+j},\\
 \Omega(Q_i,Q_{L+j})&=\frac{u_{ij}}{\sqrt2\bar c_i\bar c_{L+j}},&
 \Omega(Q_{L+j},Q_i)&=-\frac{\bar u_{ij}}{\sqrt2c_ic_{L+j}}.
\end{aligned}
\end{equation}
For example $r_i\sim-\nu/c_i$ at $Q_i$, and the last eigenvector
rows give $f_k^{\trans}r=w_{ki}/c_i$. The opposite outer block then
gives the kernel recovery, with slope difference two. At an outer
point, in coordinate $1/\nu$, the relevant $h$ has principal
coefficient $-2$ and $d\Omega_0$ has coefficient $-\kappa_a$, so
the recovered weights agree with \cref{eq:general-residue-weights}.
Thus \cref{eq:multipoint-mode-arrays} recovers every mode, with
$e_k=1$. Weighted reference anchors give other generic charts.
The forward theorem and uniqueness of the polynomial ODE identify
the evolution. Restore \cref{eq:canonical-phase-gauge} for $H$.

For repeated phases apply this argument in the constant active bases
of Appendix~\ref{app:BA-active-converse}. Their resolved infinity
directions supply the same normalized leading values. Restore
inactive zero directions and invert the canonical bases afterwards.
\end{proof}

For comparison with a prescribed realization, its fixed $e_k$ gauges
are set to one by the puncture-vertex canonical gauge transformations.
The common rescaling of \cref{sec:general-BA-interpolation} accounts
for the single redundancy among all vertex normalization constants.
The reference partition of unity in \cref{sec:multipoint-BA-data}
then makes \cref{eq:BA-converse-constants} recover the original
weights and gauges directly.
